\subsection{Partiality}
\label{sec:partiality}

The comparisons \pyinline{==}, \pyinline{!=}, \pyinline{in} and \pyinline{not in} are given in terms of
structural equality $\eqOp$ and membership $\contains$. When defined, the
operators agree with Python, yielding $\aborts{\kappa}$ when Python raises an
exception. When undefined\paperonly{ (\secref{no-coercions})}, evaluation of the expression has no derivation, and an
implementation may abort or produce the result that Python gives (\paperonly{\secref{python-compatibility}}\speconly{\secref{conformance}}). Because
equality is partial, the order in which elements are compared is observable: equality and membership short-circuit
at the first element that decides the result, so an undefined comparison later in a sequence does not make
the whole comparison undefined.

For a run that reaches an undefined operation the specification prescribes nothing: a conforming
implementation may produce Python's result, as a stock interpreter does, or abort
(\secref{python-compatibility}).

Structural equality $\eqOp$ compares literals by value, lists and tuples elementwise, dictionaries with
the same keys by their values, and objects of the same class by their fields; \pyinline{None} compares
with any value. Otherwise equality is undefined: between values of different kinds, such as a number and a
string, a list and a tuple, or a Boolean and a number; between functions; and between two
\pyinline{nan}s. Python gives an answer in each case (\pyinline{False} for unrelated kinds, identity for
functions, and \pyinline{True} for a list holding a \pyinline{nan} compared with itself), so each is a
dynamically excluded comparison (\secref{no-coercions}). Comparison stops at the first element that
decides the result, so its order is observable:

\begin{python}
print([1, "a"] == [2, 3])  # False, decided at the first element
print([1, "a"] == [1, 3])  # undefined: "a" against 3
\end{python}

\noindent Dictionaries are compared in the entry order of the left operand for the same reason. Membership
$\memOp$ is defined through equality and stops in the same way. The strict operators are given by
$\hat{\odot}$ and $\hat{\ominus}$ (\figref{operators}); the definitions of all of these are given in
\specification.

\begin{figure}
\small
\begin{center}
\input{spec/functions/binop-meaning}

\vspace{2mm}

\begin{tabular}{>{\raggedright\arraybackslash}p{0.22\textwidth}>{\raggedright\arraybackslash}p{0.62\textwidth}}
\arrayrulecolor{gray!60}\hline
\rowcolor{grammarheadcolor}$\ominus$ & $\hat{\ominus}(v)$ \\
\pyinline{not} & $\val{\pyinline{True}}$ if $v = \pyinline{False}$, and $\val{\pyinline{False}}$ if $v = \pyinline{True}$ \\
\pyinline{+} & $\val{v}$ if $v$ is a number, and $\aborts{\errType}$ otherwise \\
\pyinline{-} & $\val{-v}$ if $v$ is a number, and $\aborts{\errType}$ otherwise
\end{tabular}

\end{center}
\caption{Operator meanings: value, abort, or undefined if no case applies}
\label{fig:operators}
\end{figure}

\begin{figure}
\small
\begin{align*}
\memOp(\exList{\seq{v}}, v) &{}= \memElems(\seq{v}, v) \\*
\memOp(\exTuple{\seq{v}}, v) &{}= \memElems(\seq{v}, v) \\*
\\
\memOp(\delta, w) &{}= \pyinline{True} && \text{if } w \in \dom(\delta) \\*
\memOp(\delta, w) &{}= \pyinline{False} && \text{if } w \notin \dom(\delta) \\*
\\
\memOp(w', w) &{}= \pyinline{True} && \text{if } w' = w_1 \concat w \concat w_2 \text{ for some } w_1, w_2 \\*
\memOp(w', w) &{}= \pyinline{False} && \text{otherwise}
 \\
\\
\containsElems(\varepsilon, v) &{}= \pyinline{False} \\*
\containsElems(v' \cons \seq{v}, v) &{}= \pyinline{True} && \text{if } \eqOp(v, v') = \pyinline{True} \\*
\containsElems(v' \cons \seq{v}, v) &{}= \containsElems(\seq{v}, v) && \text{if } \eqOp(v, v') = \pyinline{False}

\end{align*}
\caption{Membership, undefined where no case applies}
\label{fig:membership}
\end{figure}


The remaining metafunctions the rules use are collected in \figref{eval-aux}: $\outcome$ reads a
block's result as the outcome of calling it, $\iterOp$ gives the sequence a generator draws from, $\getitem$
interprets subscripting, $\dispatch$ selects the case a $\pyinline{match}$ takes, and $\entries$ and
$\update$ build a dictionary from its entries.

\begin{figure}
\small
\begin{align*}
\dispatch(\rho, v, \varepsilon, \varepsilon) &= (\varnothing, \exPass) \\*
\dispatch(\rho, v, p \cons \seq{p}, s \cons \seq{s}) &= (\rho', s) && \text{if } \rho, p, v \matches \matchOk{}\rho' \\*
\dispatch(\rho, v, p \cons \seq{p}, s \cons \seq{s}) &= \dispatch(\rho, v, \seq{p}, \seq{b}) && \text{if } \rho, p, v \matches \noMatch
 \\
\\
\entries(\delta, \varepsilon) &{}= \delta \\*
\entries(\delta, \bind{w}{v} \cons \delta') &{}= \entries(\update(\delta, w, v), \delta')
 \\
\\
\update(\delta, \varepsilon) &{}= \delta \\*
\update(\delta, \bind{w}{v} \cons \delta') &{}= \update(\setitem(\delta, w, v), \delta')

\end{align*}
\caption{Metafunctions used by statement and expression evaluation}
\label{fig:eval-aux}
\end{figure}


A qualified name is resolved through the environments of loaded modules: $\env(\rho, q)$ is the
environment of module $q$ as loaded in $\rho$, and $\resolveClass(\rho, q)$ the class entry that $q$
denotes (\figref{resolution}).

\begin{figure}
\small
\begin{align*}
\env(\rho, x) &= \rho' && \text{where } \rho(x) = \modLoaded{q'}{\rho'} \\*
\env(\rho, q.x) &= \rho' && \text{where } \env(\rho, q)(x) = \modLoaded{q'}{\rho'}
 \\
\\
\resolveClass(\rho, c) &= C && \text{if } \rho(c) = \Type{C} \\*
\resolveClass(\rho, q.c) &= C && \text{if } \env(\rho, q)(c) = \Type{C}

\end{align*}
\caption{Module and class resolution}
\label{fig:resolution}
\end{figure}

